\centerline{\bf \LARGE Рыцари, лжецы и другие...}

\begin{flushright}
        \scriptsize{Цитата: 20:31 прибыл Годжо Сатору... и решил все задачи}
\end{flushright}

\problem{1} 
Однажды в четверг после дождя между аборигенами Тимом и Томом произошел следующий диалог: "Ты можешь сказать, что я рыцарь" --- гордо заявил Тим. "Ты можешь сказать, что я лжец" --- грустно ответил ему Том. Кем являются Тим и Том?

\problem{2} 
Некогда перед судом предстали три аборигена, которых для конфиденциальности мы обозначим А, Б и В. Известно, что преступление совершил ровно один из них, но кто из них является рыцарями, а кто --- лжецами, было неизвестно. "Б лжец. Но преступление совершил В" --- заявил А. "А и В либо оба рыцари, либо оба лжецы" --- сообщил суду Б. "Б говорит правду. Но тем не менее он и совершил преступление" --- сказал В. Поразмышляв недолго, судья не только сумел определить, кто есть кто, но и изобличить преступника. А вы сумеете это сделать?


\problem{3} 
В клетках квадрата 4x4 стоят аборигены. В некоторый момент каждый из них произнес: «Во всех соседних со мной клетках стоят лжецы». Какое наибольшее количество лжецов могло быть среди них?


\problem{4}
В ряду стоят четыре аборигена: два лжеца и два рыцаря. Один из них перевёртыш. Каждого спросили, не является ли он перевёртышем. Трое ответили "да", а один "нет". Кто из них является перевёртышем, если лжецы не стоят рядом?

\problem{5} 
Аборигены, среди которых есть хотя бы один рыцарь и хотя бы один лжец, встали в круг. Каждый из них сказал: «Мой сосед справа – перевёртыш». Сколько рыцарей встало в круг, если известно, что всего аборигенов 31, а перевёртышей среди них 13?

\problem{6}
В каждой клетке квадрата 3×3 стоит по аборигену. Сначала каждый из них сказал: «Каждый из моих соседей говорит правду». Затем часть из них (хотя бы один) сказала: «Все мои соседи – рыцари», а остальные (хотя бы один) сказали: «Все мои соседи – лжецы». Соседями считаются соседи по стороне. Какое наименьшее число перевёртышей может быть?

\problem{7}
В ряд стоит 31 лжец. Сколько среди них перевёртышей, если известно, что в одной расстановке все сказали правду, а в другой – все солгали?

\problem{8}
В круг встали 18 аборигенов. Известно, что среди них ровно один перевёртыш. Путешественник ОлЮр может задать каждому из аборигенов только один вопрос: «Ты – перевёртыш?». Сможет ли он найти перевёртыша? А если все аборигены – рыцари?


\hspace{3mm}

\centerline{\bf \Huge Памятка}

\problem{1} На острове живут аборигены двух племён: рыцари, которые в обычном состоянии обладают способностью всегда говорить правду, и лжецы, которые в обычном состоянии обладают способностью всегда говорить ложь.

\problem{2} Среди представителей обоих племён встречаются перевёртыши, которые воздействуют на своих соседей и изменяют их способности на противоположные.

\problem{3} Рыцари говорят правду, если не находятся под влиянием одного соседа, который перевёртыш, и его действие не нейтрализовано действием другого соседа-перевёртыша. В противном случае они лгут.

\problem{4} Лжецы лгут, если не находятся под влиянием одного соседа, который перевёртыш, и его действие не нейтрализовано действием другого соседаперевёртыша. В противном случае они говорят правду.

\problem{5} Перевёртыши воздействуют друг на друга по тем же правилам, которым подчиняются остальные члены их племён.

\problem{6} Если перевёртыши воздействуют с двух сторон (с чётного числа сторон) одновременно на аборигена, то их действия нейтрализуются.

\problem{7} Если перевёртыши воздействуют с одной стороны (с нечётного числа сторон) одновременно на аборигена, то его способности меняются на противоположные.

\problem{0.1} Давным давно абориген Гигачад сказал своим друзьям: - Вчера мой сосед заявил мне, что он лжец! Кем является Гигачад — рыцарем или лжецом?

\problem{0.2}
Как-то раз встретились два аборигена и один сказал другому: «По крайней мере один из нас – лжец». История умалчивает, ответил ли ему на это что-либо собеседник. Тем не менее определите, кем являются оба.

\problem{0.3}
В другой раз встретились два аборигена Сигма и Шкибиди. "По крайней мере один из нас – рыцарь" - глубокомысленно изрек Сигма. "Но ты то уж точно лжец!" - рассмеялся ему в лицо Шкибиди. Определите, кем являются оба.

